\newpage
\section{Implementacja}		%4

\subsection{Omówienie wybranych fragmentów kodu źródłowego}
Implementacja projektu została podzielona na modularne pliki nagłówkowe oraz pliki źródłowe, co zwiększa czytelność kodu oraz ułatwia jego konserwację. Poniżej przedstawiono kluczowe fragmenty oprogramowania realizujące najważniejsze założenia algorytmiczne i architektoniczne.

\subsubsection{Bezpieczne wstawianie elementu pod wskazany indeks}
Jedną z bardziej złożonych operacji w liście dwukierunkowej jest wstawienie elementu w środek struktury. Wymaga to precyzyjnego przepięcia wskaźników zarówno dla nowego węzła, jak i jego bezpośrednich sąsiadów:

\begin{lstlisting}[caption=Metoda addAtIndex, label={lst:listing-cpp}, language=C++]
void LinkedList::addAtIndex(int index, int value) {
    if (index < 0 || index > size) {
        std::cout << "Nieprawidlowy indeks!\n";
        return;
    }
    if (index == 0) { addFront(value); return; }
    if (index == size) { addBack(value); return; }

    Node* current = head;
    for (int i = 0; i < index; i++) {
        current = current->next;
    }

    Node* newNode = NodeFactory::createNode(value); // Użycie wzorca Fabryki
    Node* prevNode = current->prev;

    prevNode->next = newNode;
    newNode->prev = prevNode;
    newNode->next = current;
    current->prev = newNode;

    size++;
}
\end{lstlisting}

\subsubsection{Implementacja Wzorca Iteratora}
Wzorzec Iteratora\cite{gamma_design_patterns} pozwala na sekwencyjny dostęp do elementów kolekcji bez ujawniania jej wewnętrznej budowy. Poniższy fragment prezentuje przeciążenie operatorów inkrementacji oraz wyłuskania:

\begin{lstlisting}[caption=Wzorzec projektowy - iterator, label={lst:listing-cpp}, language=C++]
class ListIterator {
private:
    Node* current;
public:
    explicit ListIterator(Node* ptr) : current(ptr) {}

    int& operator*() const {
        return current->data;
    }

    ListIterator& operator++() {
        if (current != nullptr) {
            current = current->next;
        }
        return *this;
    }

    bool operator!=(const ListIterator& other) const {
        return current != other.current;
    }
};
\end{lstlisting}

\subsection{Opis uzyskanych wyników i działania programu}
W wyniku poprawnej kompilacji i uruchomienia programu (zarówno za pomocą skryptu konsolowego MSVC \texttt{cl.exe}, jak i przez środowisko Visual Studio) otrzymano w pełni stabilną aplikację konsolową. Uzyskane rezultaty obejmują:
\begin{enumerate}
    \item \textbf{Dynamiczne zarządzanie pamięcią:} Program poprawnie alokuje pamięć dla nowych węzłów w stercie oraz – dzięki zaimplementowanemu destruktorowi i metodzie \texttt{clear()} – zwalnia ją całkowicie w momencie zamknięcia, nie powodując wycieków pamięci (\textit{memory leaks}).
    \item \textbf{Poprawność strukturalna:} Operacje wstawiania oraz usuwania elementów na początku, końcu oraz pod wskazanym indeksem zachowują spójność wskaźników \texttt{next} oraz \texttt{prev}.
    \item \textbf{Efektywność iteracji:} Zastosowanie dedykowanego iteratora umożliwiło czytelne i bezpieczne przeglądanie zawartości struktury danych, co zbliża implementację do standardów biblioteki STL języka C++.
    \item \textbf{Generowanie dokumentacji:} Użycie komentarzy zgodnych ze standardem Doxygen pozwoliło na automatyczne wyczyszczenie struktury kodu i przygotowanie kompletnego raportu API projektu.
\end{enumerate}